\documentclass[12pt]{article}

\usepackage{amsmath}
\usepackage{amsthm}
\usepackage{thmtools}
\usepackage{amssymb}
\usepackage{enumitem}
\usepackage{float}
\usepackage{gensymb}
\usepackage{geometry}
\usepackage{graphicx}
\usepackage{hyperref}
\usepackage{mathtools}
\usepackage{multirow}
\usepackage{physics}
\usepackage{tabularx}
\usepackage{wrapfig}
\usepackage[utf8]{inputenc}


\setlist{parsep=1em}

\newcommand{\succr}[1]{#1 \! +\!\!\!+}

\newenvironment{df}[2][\relax]{
    \noindent\setlength{\parskip}{1em}%
    \ifx#1\relax%
        \textbf{#2. }%
    \else%
        \textbf{#2} (#1)\textbf{. }%
    \fi%
    \ignorespaces%
}{\vspace{4mm}}

\theoremstyle{remark}
\newtheorem*{remark}{Remark}


\title{Definitions in Analysis I by Terence Tao}
\author{Li Thon Hong}
\date{May 2026}

\begin{document}

\maketitle
\tableofcontents
\vspace{5mm}
\begin{remark}
    This document follows the third edition of \textit{Analysis I}.
\end{remark}

\pagebreak


\section{Introduction}
No definitions found in this chapter.

\pagebreak


\section{Starting at the beginning: the natural numbers}
\subsection{The Peano axioms}

\begin{df}{Axiom 2.1}
    $0$ is a natural number.
\end{df}


\begin{df}{Axiom 2.2}
    If $n$ is a natural number, then $\succr{n}$ is also a natural number.
\end{df}


\begin{df}{Definition 2.1.3}
    We define $1$ to be the number $\succr{0}$, $2$ to be the number $\succr{(\succr{0})}$, 3 to be the number $\succr{(\succr{(\succr{0})})}$, etc.
    In other words, $1 \coloneqq \succr{0}$, $2 \coloneqq \succr{1}$, $3 \coloneqq \succr{2}$, etc.
\end{df}


\begin{df}{Axiom 2.3}
    $0$ is not the successor of any natural number, i.e., we have $\succr{n} \ne 0$ for every natural number $n$.
\end{df}


\begin{df}{Axiom 2.4}
    Different natural numbers must have different successors; i.e. if $n$, $m$ are natural numbers and $n \ne m$, then $\succr{n} = \succr{m}$.
    Equivalently, if $\succr{n} = \succr{m}$, then we must have $n = m$.    
\end{df}


\begin{df}[Principle of mathematical induction]{Axiom 2.5}
    Let $P(n)$ be any property pertaining to a natural number $n$.
    Suppose that $P(0)$ is true, and suppose that whenever $P(n)$ is true, $P(\succr{n})$ is also true.
    Then $P(n)$ is true for every natural number $n$.
\end{df}


\subsection{Addition}

\begin{df}[Addition of natural numbers]{Definition 2.2.1}
    Let $m$ be a natural number.
    To add $0$ to $m$, we define $0 + m \coloneqq m$.
    Now, suppose inductively that we have defined how to add $n$ to $m$.
    Then, we can add $\succr{n}$ to m by defining $(\succr{n}) + m \coloneqq \succr{(n + m)}$.%
\end{df}


\begin{df}{Lemma 2.2.2}
    For any natural number $n$, $n + 0 = n$.
\end{df}


\begin{df}{Lemma 2.2.3}
    For any natural numbers $n$ and $m$, $n + (\succr{m}) = \succr{(n + m)}$.
\end{df}


\begin{df}[Addition is commutative]{Proposition 2.2.4}
    For any natural numbers $n$ and $m$, $n + m = m + n$.
\end{df}


\begin{df}[Addition is associative]{Proposition 2.2.5}
    For any natural numbers $a$, $b$, $c$, we have $(a + b) + c = a +(b + c)$.
\end{df}


\begin{df}[Cancellation law]{Proposition 2.2.6}
    Let $a$, $b$, $c$ be natural numbers such that $a + b = a + c$.
    Then we have $b = c$.
\end{df}


\begin{df}[Positive natural numbers]{Definition 2.2.7}
    A natural number $b$ is said to be \emph{positive} iff it is not equal to $0$.    
\end{df}


\begin{df}{Proposition 2.2.8}
    If $a$ is positive and $b$ is a natural number, then $a + b$ is positive.
\end{df}


\begin{df}{Corollary 2.2.9}
    If $a$ and $b$ are natural numbers such that $a + b = 0$, then $a = 0$ and $b = 0$.
\end{df}


\begin{df}{Lemma 2.2.10}
    Let $a$ be a positive number.
    Then there exists exactly one natural number $b$ such that $\succr{b} = a$.
\end{df}


\begin{df}[Ordering of the natural numbers]{Definition 2.2.11}
    Let $n$ and $m$ be natural numbers.
    We say that $n$ is greater than or equal to $m$, and write $n \ge m$ or $m \le n$, iff we have $n = m + a$ for some natural number $a$.
    We say that $n$ is \emph{strictly greater than} $m$, and write $n > m$ or $m < n$, iff $n \ge m$ and $n \ne m$.
\end{df}


\begin{df}[Basic properties of order for natural numbers]{Proposition 2.2.12}
    Let $a$, $b$, $c$ be natural numbers. Then
    \begin{enumerate}[topsep=0mm,parsep=0mm,label=(\alph*)]
        \item \textit{Reflexivity: $a \ge a$.}
        \item \textit{Transitivity: If  $a \ge b$ and $b \ge c$, then $a \ge c$.}
        \item \textit{Anti-symmetry: If $a \ge b$ and $b \ge a$, then $a = b$.}
        \item \textit{Preservation of order: $a \ge b$ if and only if $a + c \ge b + c$.}
        \item \textit{$a < b$ if and only if $\succr{a} < b$.}
        \item \textit{$a < b$ if and only if $b = a + d$ for some positive number $d$.}
    \end{enumerate}    
\end{df}


\begin{df}[Trichotomy of order for natural number]{Proposition 2.2.13}
    Let $a$ and $b$ be natural numbers.
    Then exactly one of the following statements is true: $a < b$, $a = b$, or $a > b$.
\end{df}


\begin{df}[Strong principle of induction]{Proposition 2.2.14}
    Let $m_0$ be a natural number, and let $P(m)$ be a property pertaining to an arbitrary natural number $m$.
    Suppose that for each $m \ge m_0$, we have the following implication: if $P(m')$ is true for all natural numbers $m_0 \le m' < m$, then $P(m)$ is also true.
    Then we can conclude that $P(m)$ is true for all natural numbers $m \ge m_0$.
\end{df}


\subsection{Multiplication}

\begin{df}[Multiplication of natural numbers]{Definition 2.3.1}
    Let $m$ be a natural number.
    To multiply $0$ to $m$, we define $0 \times m \coloneqq 0$.
    Now, suppose inductively that we have defined how to multiply $n$ to $m$.
    Then, we can multiply $\succr{n}$ to m by defining $(\succr{n}) \times m \coloneqq (n \times m) + m$.%
\end{df}


\begin{df}[Multiplication is commutative]{Lemma 2.3.2}
    Let $n$, $m$ be natural numbers.
    Then $n \times m = m \times n$.    
\end{df}


\begin{df}[Positive natural numbers have no zero divisors]{Lemma 2.3.3}
    Let $n$, $m$ be natural numbers.
    Then $n \times m = 0$ if and only if at least one of $n$, $m$ is equal to $0$.
    In particular, if $n$ and $m$ are both positive, then $nm$ is also positive.
\end{df}


\begin{df}[Distributive law]{Proposition 2.3.4}
    For any natural numbers $a$, $b$, $c$, we have $a(b + c) = ab + ac$ and $(b + c)a = ba + ca$.
\end{df}


\begin{df}[Multiplication is associative]{Proposition 2.3.5}
    For any natural numbers $a$, $b$, $c$, we have $(a \times b)\times c = a \times (b \times c)$
\end{df}


\begin{df}[Multiplication preserves order]{Proposition 2.3.6}
    If $a$, $b$ are natural numbers such that $a < b$, and $c$ is positive, then $ac < bc$.
\end{df}


\begin{df}[Cancellation law]{Corollary 2.3.7}
    Let $a$, $b$, $c$ be natural numbers such that $ac = bc$ and $c$ is nonzero.
    Then $a = b$.
\end{df}


\begin{df}[Euclidean algorithm]{Proposition 2.3.9}
    Let $n$ be a natural number, and let $q$ be a positive number.
    Then there exist natural numbers $m$, $r$ such that $0 \le r < q$ and $n = mq + r$.
\end{df}


\begin{df}[Exponentiation for natural numbers]{Definition 2.3.11}
    Let $m$ be a natural number.
    To raise $m$ to the power $0$, we define $m^0 \coloneqq 1$; in particular, we define $0^0 \coloneqq 1$.
    Now suppose recursively that $m^n$ has been defined for some natural number $n$, then we define $m^{\succr{n}} \coloneqq m^n \times m$.
\end{df}


\pagebreak


\section{Set theory}
\subsection{Fundamentals}
\begin{df}[Sets are objects]{Axiom 3.1}
    If $A$ is a set, then $A$ is also an object.
    In particular, given two sets $A$ and $B$, it is meaningful to ask whether $A$ is also an element of $B$.
\end{df}


\begin{df}[Equality of sets]{Definition 3.1.4}
    Two sets $A$ and $B$ are \emph{equal}, $A = B$, iff every element of $A$ is an element of $B$ and vice versa.
    To put it another way, $A = B$ if and only if every element $x$ of $A$ belongs also to $B$, and every element $y$ of $B$ belongs also to $A$.
\end{df}


\begin{df}[Empty set]{Axiom 3.2}
    There exists a set $\varnothing$, known as the \emph{empty set}, which contains no elements, i.e., for every object $x$ we have $x \notin \varnothing$.
\end{df}


\begin{df}[Single choice]{Lemma 3.1.6}
    Let $A$ be a non-empty set.
    Then there exists an object $x$ such that $x \in A$.
\end{df}


\begin{df}[Singleton sets and pair sets]{Axiom 3.3}
    If $a$ is an object, then there exists a set $\{a\}$ whose only element is $a$, i.e., for every object $y$, we have $y \in \{ a \}$ if and only if $y = a$; we refer to $\{ a \}$ as the singleton set whose element is $a$.
    Furthermore, if $a$ and $b$ are objects, then there exists a set $\{a, b\}$ whose only elements are $a$ and $b$; i.e., for every object $y$, we have $y \in \{a, b\}$ if and only if $y = a$ or $y = b$; we refer to this set as the pair set formed by $a$ and $b$.
\end{df}


\begin{df}[Pairwise union]{Axiom 3.4}
    Given any two sets $A$, $B$, there exists a set $A \cup B$, called the \emph{union} $A \cup B$ of $A$ and $B$, whose elements consist of all the elements which belong to $A$, or $B$ or both.
    In other words, for any object $x$,
    {\setlength{\belowdisplayskip}{-10pt} \setlength{\belowdisplayshortskip}{-10pt}
    \[
        x \in A \cup B \iff (x \in A \text{ or } x \in B).
    \]
    }
\end{df}


\begin{df}{Lemma 3.1.13}
    If $a$ and $b$ are objects, then $\{ a, b \} = \{ a \} \cup \{ b \}$.
    If $A$, $B$, $C$ are sets, then the union operation is commutative (i.e., $A \cup B = B \cup A$) and associative (i.e., $(A \cup B) \cup C = A \cup (B \cup C)$).
    Also, we have $A \cup A = A \cup \varnothing = \varnothing \cup A = A$.
\end{df}


\begin{df}[Subsets]{Definition 3.1.15}
    Let $A$, $B$ be sets.
    We say that $A$ is a \emph{subset} of $B$, denoted $A \subseteq B$, iff every element of $A$ is also an element of $B$, i.e.
    \[
        \text{For any object } x, \quad x \in a \implies x \in B.
    \]
    We say that $A$ is a \emph{proper subset} of $B$, denoted $A \subsetneq B$, if $A \subseteq B$ and $A \ne B$.
\end{df}


\begin{df}[Sets are partially ordered by set inclusion]{Proposition 3.1.18}
    Let $A$, $B$, $C$ be sets.
    If $A \subseteq B$ and $B \subseteq C$, then $A \subseteq C$.
    If $A \subseteq B$ and $B \subseteq A$, then $A = B$.
    Finally, if $A \subsetneq B$ and $B \subsetneq C$, then $A \subsetneq C$.
\end{df}


\begin{df}[Axiom of specification/axiom of specification]{Axiom 3.5}
    Let $A$ be a set, and for each $x \in A$, let $P(x)$ be a property pertaining to $x$ (i.e., $P(x)$ is either a true statement or a false statement).
    Then there exists a set, called $\{x \in A : P(x) \text{ is true} \}$ (or simply $\{x \in A : P(x)\}$ for short), whose elements are precisely the elements $x$ in $A$ for which $P(x)$ is true.
    In other words, for any object $y$,
    {\setlength{\belowdisplayskip}{-10pt} \setlength{\belowdisplayshortskip}{-10pt}
    \[
        y \in \{ x \in A : P(x) \text{ is true} \} \iff (y \in A \text{ and } P(y) \text{ is true}).
    \]
    }
\end{df}


\begin{df}[Intersections]{Definition 3.1.23}
    The \emph{intersection} $S_1 \cap S_2$ of two sets is defined to be the set
    {\setlength{\belowdisplayskip}{-10pt} \setlength{\belowdisplayshortskip}{-10pt}
    \[
        S_1 \cap S_2 \coloneqq \{ x \in S_1 : x \in S_2 \}.
    \]
    }
\end{df}

\begin{df}[Difference set]{Definition 3.1.27}
    Given two sets $A$ and $B$, we define the set $A - B$ or $A \setminus B$ to be the set $A$ with any elements of $B$ removed:
    {\setlength{\belowdisplayskip}{-10pt} \setlength{\belowdisplayshortskip}{-10pt}
    \[
        A \setminus B \coloneqq \{ x \in A : x \notin B \}.
    \]
    }
\end{df}


\begin{df}[Sets form a boolean algebra]{Proposition 3.1.28}
    Let $A, B, C$ be sets, and let $X$ be a set containing $A, B, C$ as subsets.
    \begin{enumerate}[topsep=0mm,parsep=0mm,label=(\alph*)]
        \item (Minimal element) We have $A \cup \varnothing = A$ and $A \cap \varnothing = \varnothing$.
        \item (Maximal element) We have $A \cup X = X$ and $A \cap X = A$.
        \item (Identitiy) We have $A \cap A= A$ and $A \cup A = A$.
        \item (Commutativity) We have $A \cup B = B \cup A$ and $A \cap B = B \cap A$.
        \item (Associativity) We have $(A \cup B) \cup C = A \cup (B \cup C)$and $(A \cap B) \cap C = A \cap (B \cap C)$.
        \item (Distributivity) We have $A \cap (B \cup C) = (A \cap B) \cup (A \cap C)$ and $A \cup (B \cap C)=(A \cup B) \cap (A \cup C)$.
        \item (Partition) We have $A \cup (X \setminus A) = X$ and $A \cap (X \setminus A) = \varnothing$.
        \item (De Morgan Laws) We have $X \setminus (A \cup B) = (X \setminus A) \cap (X \setminus B)$ and $X \setminus (A \cap B) = (X \setminus A) \cup (X \setminus B)$.
    \end{enumerate}
\end{df}


\begin{df}[Replacement]{Axiom 3.6}
    Let $A$ be a set.
    For any object $x \in A$ and any object $y$, suppose we have a statement $P(x, y)$ pertaining to $x$ and $y$, such that for each $x \in A$ there is at most one $y$ for which $P(x, y)$ is true.
    Then there exists a set $\{ y : P(x, y) \text{ is true for some } x \in A \}$, such that for any object $z$,
    {\setlength{\belowdisplayskip}{-10pt} \setlength{\belowdisplayshortskip}{-10pt}
    \[
        z \in \{ y : P(x, y) \text{ is true for some } x \in A \} \iff P(x, z) \text{ is true for some } x \in A.
    \]
    }
\end{df}


\begin{df}[Infinity]{Axiom 3.7}
    There exists a set $\mathbb{N}$, whose elements are called natural numbers, as well as an object $0$ in $\mathbb{N}$, and an object $\succr{n}$ assigned to every natural number $n \in \mathbb{N}$, such that the Peano axioms (Axioms 2.1 - 2.5) hold.
\end{df}


\subsection{Russell's paradox}
\begin{df}[Universal specification]{Axiom 3.8}
    (\textit{Dangerous!})
    Suppose for every object $x$ we have a property $P(x)$ pertaining to $x$ (so that for every $x$, $P(x)$ is either a true statement or a false statement).
    Then there exists a set $\{x : P(x) \text{ is true} \}$ such that for every object $y$,
    {\setlength{\belowdisplayskip}{-10pt} \setlength{\belowdisplayshortskip}{-10pt}
    \[
        y \in \{x : P(x) \text{ is true} \} \iff P(y) \text{ is true.}
    \]
    }
\end{df}


\begin{df}[Regularity]{Axiom 3.9}
    If $A$ is a non-empty set, then there is at least one element $x$ of $A$ which is either not a set, or is disjoint from $A$.
\end{df}

\end{document}